\section{Related Work}
\subsection{Foundation Models for Visual Understanding}
Visual understanding requires both semantic recognition and spatial localization. Vision Transformers represent images as patch tokens~\cite{vit}, while self-supervised learning in DINO produces features that capture object structure without task-specific labels~\cite{dino}. DINOv2 and DINOv3 extend this representation approach at scale~\cite{dinov2,dinov3}. Their dense features support correspondence across images, providing a basis for transferring an annotated target from a reference to a query.

Promptable foundation models expose spatial and semantic information through a segmentation interface. SAM predicts masks from spatial prompts~\cite{sam}; SAM~2 extends this interface to video~\cite{sam2}; and SAM~3 supports concept-level prompts~\cite{sam3}. CLIP connects visual representations with language through image--text alignment~\cite{clip}, while models such as Qwen3-VL support multimodal reasoning~\cite{qwen3vl}. Together, these capabilities enable the two routes studied here: identifying a target through dense visual correspondence, or describing it semantically and grounding that description with a promptable segmenter.

\subsection{In-Context Segmentation}
Reference-conditioned segmentation builds on one-shot prediction~\cite{oslsm}, support prototypes~\cite{panet}, and dense support--query correspondence~\cite{hsnet}. Visual in-context learning broadens this idea by expressing tasks through image examples: visual prompting casts prediction as image completion~\cite{visualprompt}, Painter learns across dense prediction tasks~\cite{painter}, and SegGPT provides a general context-conditioned segmentation interface~\cite{seggpt}. These approaches establish the annotated example as a flexible specification of the desired output.

Recent methods differ in how they turn that example into a query mask. SegIC turns foundation-model correspondences into geometric, visual, and meta instructions~\cite{segic}; SINE introduces in-context interaction~\cite{sine}; and VRP-SAM learns a reference-prompt encoder~\cite{vrpsam}. Frozen-model pipelines instead reuse pretrained representations and segmentation priors: PerSAM derives target-specific information from one example~\cite{persam}, and Matcher links general feature matching to promptable segmentation~\cite{matcher}. INSID3 operates directly on DINOv3 features~\cite{insid3}, whereas CG-ICS combines generated concepts, search, and visual exemplars~\cite{cgics}. We compare these last two complete inference systems and then examine their individual processing stages.

The chosen reference is itself a source of variation. Good Examples studies example selection~\cite{goodexamples}, UNICL-SAM addresses uncertainty in imperfect examples~\cite{unicl}, and CG-ICS explicitly targets reference robustness~\cite{cgics}. Our reference study measures performance when the query is fixed and the supplied demonstration changes. Cross-domain work provides a complementary perspective: RePRI uses transductive optimization~\cite{repri}, while ABCDFSS adapts feature-processing layers at test time~\cite{abcdfss}. Their adaptation mechanisms differ from the frozen inference evaluated here, motivating a separate domain-transfer analysis.
